The first rank where divisibility stops characterizing the
series-parallel types is $n=5$.  We treat $S_5$ as the first nontrivial
test of the extension-tiler converse.

We first isolate the finite classification input.  Let $N$ be the
four-element poset on $\{a,b,c,d\}$ with relations
\[
        a<b,\qquad c<b,\qquad c<d.
\]
This is the forbidden induced poset in the Valdes--Tarjan--Lawler
characterization of finite series-parallel posets, viewing a poset through
its transitive directed comparability graph.

\begin{lemma}[One-point $N$-extensions]
\label{lem:one-point-n-extensions}
Every non-series-parallel poset on five elements is isomorphic to a
one-point extension of $N$.

More explicitly, let $x$ be the added point and put
\[
        D=\{y\in N:y<x\},\qquad U=\{y\in N:x<y\}.
\]
Then $D$ is an order ideal of $N$, $U$ is an order filter of $N$, and
\[
        d<_N u\quad\text{for every }d\in D,\ u\in U.
\]
Conversely, any pair $(D,U)$ satisfying these three conditions defines a
five-element poset containing $N$ as an induced subposet.
\end{lemma}

\begin{proof}
By the Valdes--Tarjan--Lawler characterization, a finite poset is
non-series-parallel exactly when it contains an induced $N$-poset.  Thus a
non-series-parallel poset on five elements is obtained from an induced
copy of $N$ by adding one further point $x$.

If $y<x$ and $z<y$ in $N$, then $z<x$, so $D$ is an ideal.  Dually, if
$x<y$ and $y<z$ in $N$, then $x<z$, so $U$ is a filter.  Finally, if
$d\in D$ and $u\in U$, then $d<x<u$, hence $d<u$ by transitivity.  Since
the chosen four-point subposet is induced, this comparability must already
belong to $N$.

Conversely, if $D$ is an ideal, $U$ is a filter, and every $d\in D$ is below
every $u\in U$ in $N$, then adjoining relations $d<x$ for $d\in D$ and
$x<u$ for $u\in U$ creates no new comparability among the four old
elements beyond those already in $N$.  Hence $N$ remains an induced
subposet.
\end{proof}

\begin{lemma}[Profile table]
\label{lem:s5-profile-table}
Up to isomorphism, there are exactly $15$ non-series-parallel posets on
five elements.  Their numbers of linear extensions are
\[
5,5,7,7,8,9,9,11,14,14,16,16,18,18,25.
\]
\end{lemma}

\begin{proof}
The ideals of $N$ are
\[
\emptyset,\ a,\ c,\ ac,\ cd,\ abc,\ acd,\ abcd,
\]
and the filters of $N$ are
\[
\emptyset,\ b,\ d,\ ab,\ bd,\ abd,\ bcd,\ abcd.
\]
Here, for instance, $ac$ denotes the set $\{a,c\}$.

The valid profiles $(D,U)$ from Lemma~\ref{lem:one-point-n-extensions}
are exactly the nonempty entries of the following matrix.  Each entry is
the value of $|L(P_{D,U})|$ for the corresponding one-point extension:
\[
\begin{array}{c|cccccccc}
D\backslash U
&\emptyset& b& d& ab& bd& abd& bcd& abcd\\
\hline
\emptyset &25&18&16&9&14&8&7&5\\
a         &16&9&&&&&&\\
c         &18&11&9&&7&&&\\
ac        &14&7&&&&&&\\
cd        &9&&&&&&&\\
abc       &7&&&&&&&\\
acd       &8&&&&&&&\\
abcd      &5&&&&&&&
\end{array}
\]

The entries are computed by insertion into the five linear extensions of
$N$:
\[
        acbd,\quad acdb,\quad cabd,\quad cadb,\quad cdab.
\]
For a profile $(D,U)$, the number $|L(P_{D,U})|$ is the sum, over these
five words, of the number of slots in which $x$ may be inserted after all
elements of $D$ and before all elements of $U$.

There are $20$ valid profiles before quotienting by isomorphism.  The only
identifications among them are
\[
\begin{gathered}
(ac,b)\cong(abc,\emptyset),\qquad
(\emptyset,bcd)\cong(c,bd),\\
(\emptyset,abd)\cong(acd,\emptyset),\qquad
(c,d)\cong(cd,\emptyset),\\
(\emptyset,ab)\cong(a,b).
\end{gathered}
\]
One may check these five equivalences by the following explicit relabelings
of $\{a,b,c,d,x\}$, written in cycle notation:
\[
\begin{array}{c|c}
\text{profiles} & \text{relabeling}\\
\hline
(ac,b)\cong(abc,\emptyset) & (b\,x)\\
(\emptyset,bcd)\cong(c,bd) & (c\,x)\\
(\emptyset,abd)\cong(acd,\emptyset) & (a\,d\,b\,x\,c)\\
(c,d)\cong(cd,\emptyset) & (d\,x)\\
(\emptyset,ab)\cong(a,b) & (a\,x).
\end{array}
\]

Appendix~\ref{app:profile-quotient} records a short invariant check that
there are no further identifications: for each of the $15$ quotient
classes, the pair consisting of $|L|$ and the multiset of lower/upper
comparability counts of its elements is distinct.

Thus the $20$ profiles give exactly $15$ isomorphism classes.  Reading one
representative from each class gives the displayed multiset of extension
numbers.
\end{proof}

\begin{corollary}[Divisibility reduction]
\label{cor:s5-divisibility-reduction}
Among non-series-parallel posets on five elements, only three isomorphism
classes satisfy $|L(P)|\mid 120$: the two classes with $|L(P)|=5$ and the
single class with $|L(P)|=8$.
\end{corollary}

\begin{lemma}[Fixed end-position obstruction]
\label{lem:fixed-position-obstruction}
Let $P$ be a poset on five elements with $|L(P)|=5$.  If all linear
extensions of $P$ have the same first letter, or all have the same last
letter, then $P$ is not an extension tiler.
\end{lemma}

\begin{proof}
Suppose first that every word in $L(P)$ begins with $p$.  For any
$a\in S_5$, every word in $aL(P)$ begins with $a(p)$, so the translate is
contained in one first-letter fiber
\[
        F_x^{\mathrm{first}}=\{\sigma\in S_5:\sigma_1=x\}.
\]
Each such fiber has size $4!=24$.  A tiling by translates of $L(P)$ would
partition every first-letter fiber into blocks of size $5$, impossible
because $5\nmid 24$.

The last-letter case is identical, using the fibers
$F_x^{\mathrm{last}}=\{\sigma\in S_5:\sigma_5=x\}$.
\end{proof}

\begin{lemma}[Two-letter obstruction]
\label{lem:two-letter-obstruction}
Let $P=P_{\emptyset,abd}$ be the one-point extension with relations
\[
        a<b,\qquad c<b,\qquad c<d,\qquad x<a,\qquad x<b,\qquad x<d.
\]
Then $P$ is not an extension tiler.
\end{lemma}

\begin{proof}
The linear extensions of $P$ are
\[
\begin{gathered}
cxabd,\quad cxadb,\quad cxdab,\quad xacbd,\quad xacdb,\\
xcabd,\quad xcadb,\quad xcdab.
\end{gathered}
\]
Thus their first-letter and first-two-letter distributions are
\[
\#\{\sigma\in L(P):\sigma_1=x\}=5,\qquad
\#\{\sigma\in L(P):\sigma_1=c\}=3,
\]
and
\[
\#(xc)=3,\qquad \#(cx)=3,\qquad \#(xa)=2,
\]
where $\#(uv)$ denotes the number of extensions beginning with the ordered
pair $(u,v)$.

Assume for contradiction that $S_5$ is partitioned by $15$ left
translates of $L(P)$.  For each translate $gL(P)$ write
\[
        g(x)=u,\qquad g(c)=v,\qquad g(a)=w.
\]
Then this translate contributes $3$ elements to the ordered-pair fiber
$(u,v)$, $3$ elements to the fiber $(v,u)$, and $2$ elements to the fiber
$(u,w)$.

First count only first letters.  Let
\[
\alpha_y=\#\{g:g(x)=y\},\qquad
\beta_y=\#\{g:g(c)=y\}.
\]
The first-letter fiber over $y$ has size $24$, so
\[
        5\alpha_y+3\beta_y=24.
\]
The nonnegative integer solutions are $(\alpha_y,\beta_y)=(3,3)$ and
$(0,8)$.  Since $\sum_y\alpha_y=15$, the second possibility cannot occur.
Thus
\[
        \alpha_y=\beta_y=3
\]
for every label $y$.

Now fix distinct labels $u,v$, and define
\[
N_{uv}=\#\{g:g(x)=u,\ g(c)=v\},\qquad
C_{uv}=\#\{g:g(x)=u,\ g(a)=v\}.
\]
The ordered-pair fiber $(u,v)$ has size $3!=6$, hence
\[
        3N_{uv}+3N_{vu}+2C_{uv}=6.
\]
Modulo $3$, this gives $C_{uv}\equiv 0\pmod 3$.  For fixed $u$,
\[
        \sum_{v\ne u} C_{uv}=\alpha_u=3,
\]
so there is a unique value $f(u)\ne u$ with $C_{u,f(u)}=3$, and all other
$C_{uv}$ are zero.

The fiber equation now says that if $f(u)=v$, then
$N_{uv}+N_{vu}=0$, while if $f(u)\ne v$, then $N_{uv}+N_{vu}=2$.  Applying
the same statement to the opposite ordered pair $(v,u)$ gives
\[
        f(u)=v \quad\Longleftrightarrow\quad f(v)=u.
\]
Thus $f$ would be a fixed-point-free involution on a set of five
elements.  No such involution exists.  This contradiction excludes the
tiling.
\end{proof}

\begin{theorem}[$S_5$ classification]
\label{thm:s5-classification}
For posets $P$ on five elements, $L(P)$ tiles $S_5$ by left translates if
and only if $P$ is series-parallel.
\end{theorem}

\begin{proof}
The positive direction follows from Theorem~\ref{thm:sp-extension-tiler}.
Conversely, suppose that $P$ is an extension tiler on five elements.  Then
$|L(P)|$ must divide $|S_5|=120$.

If $P$ is not series-parallel, Corollary~\ref{cor:s5-divisibility-reduction}
shows that only three non-series-parallel types pass this divisibility
test: the two types with $|L(P)|=5$ and the type with $|L(P)|=8$.

The two $|L(P)|=5$ profiles are $(abcd,\emptyset)$ and
$(\emptyset,abcd)$.  In the first, all linear extensions end with $x$; in
the second, all linear extensions begin with $x$.  Lemma~\ref{lem:fixed-position-obstruction}
excludes both types.

The remaining divisible non-series-parallel type is represented by
$P_{\emptyset,abd}$ and is excluded by
Lemma~\ref{lem:two-letter-obstruction}.  Therefore no non-series-parallel
poset on five elements is an extension tiler.
\end{proof}

For auditability, the finite table above and the former rank and packing
certificates are mirrored in the static audit file
\[
\texttt{03-certificates/s5\_expected\_output.txt}.
\]
The theorem itself uses only the finite table and the elementary
obstructions proved above.  The supplementary replay script is included
for auditability and reproduces the aggregate table together with the
former rank and packing certificates; it is not a dependency of the proof.

\begin{remark}
This $n=5$ means Coxeter type $A_4$ and the group $S_5$.  It is not the
separate geometric question of dissecting a tetrahedron into five
congruent pieces.
\end{remark}
